% !TeX root = main.tex

% Escribe cada apartado dentro de su comando. Deja una línea en blanco
% entre párrafos. Para citar, usa \cite{clave} con la clave de referencias.bib.

% Punto 5 de la guía.
\newcommand{\CarlosTrabajosRelacionados}{%
  \section{Trabajos relacionados}

  Se comparan tres antecedentes para la aplicación de gestión de laboratorios
  informáticos (CLAB) de la Universidad Técnica Estatal de Quevedo (UTEQ).
  Gil Vera et al. presentan el \emph{STEM Real Remote Laboratory} (STEM RRL)
  \cite{gilVera2024}. García-Granja et al. desarrollan Biblio-Smart para la
  Biblioteca General de la Universidad de Málaga \cite{garciaGranja2022}.
  Müller et al. estudian el modelo del \emph{International Timetabling
  Competition 2019} (ITC 2019), construido a partir de datos de universidades que
  usan UniTime \cite{muller2025}. Las funciones de CLAB se presentan como
  propuestas de diseño. Su esqueleto actual demuestra la comunicación entre
  los servicios de Laboratorios y Reservas.

  \subsection{Recurso que se reserva}

  STEM RRL permite reservar acceso remoto a equipos físicos para prácticas de
  electrónica \cite{gilVera2024}. Biblio-Smart permite reservar un asiento
  individual de biblioteca \cite{garciaGranja2022}. El modelo de ITC 2019
  asigna aulas y franjas horarias a clases durante la planificación académica
  \cite{muller2025}. CLAB propone gestionar laboratorios informáticos completos
  para actividades presenciales. De STEM RRL puede tomar la identificación
  inequívoca del recurso reservado. De ITC 2019 puede aprovechar la relación
  entre espacio, franja y actividad para que cada reserva indique qué
  laboratorio se usará, cuándo y para qué.

  \subsection{Consulta de disponibilidad}

  En STEM RRL, el usuario elige un laboratorio remoto, un día y una franja
  habilitada en el calendario \cite{gilVera2024}. Biblio-Smart distingue en un
  mapa los asientos libres, reservados y ocupados, y permite consultar sus
  características antes de elegir \cite{garciaGranja2022}. ITC 2019 representa
  los periodos en los que cada aula puede asignarse \cite{muller2025}. CLAB
  propone un calendario por laboratorio que combine la selección por fecha de
  STEM RRL con la distinción visual de estados de Biblio-Smart. Un docente
  podría ver si una franja está disponible, reservada o bloqueada por la
  institución antes de enviar la solicitud.

  \subsection{Permisos según el usuario}

  STEM RRL asigna a estudiantes, docentes y administradores permisos diferentes
  para reservar, configurar equipos o programar mantenimiento
  \cite{gilVera2024}. Biblio-Smart ofrece un flujo de reserva para los miembros
  de la universidad a través de su aplicación institucional
  \cite{garciaGranja2022}. ITC 2019 emplea matrículas y grupos para determinar
  a qué clases puede asignarse cada estudiante dentro del horario
  \cite{muller2025}. CLAB puede tomar de STEM RRL la separación entre quienes
  solicitan una franja y quienes administran la disponibilidad. Esa distinción
  permite definir quién puede solicitar una reserva, programar una serie o
  bloquear un laboratorio, de acuerdo con las responsabilidades de la UTEQ.

  \subsection{Reservas puntuales y recurrentes}

  STEM RRL permite reservar franjas por adelantado y ofrece a los docentes
  la opción de programar reservas periódicas \cite{gilVera2024}.
  Biblio-Smart describe la reserva individual de un asiento con un plazo
  para ocuparlo \cite{garciaGranja2022}. ITC 2019 admite clases que se repiten
  varios días o únicamente en determinadas semanas del semestre
  \cite{muller2025}. CLAB propone cubrir tanto una práctica puntual como una
  asignación recurrente para una materia. De STEM RRL puede tomar la
  distinción entre ambos flujos y de ITC 2019 puede aprovechar el manejo de
  excepciones, como omitir una fecha en la que el laboratorio esté bloqueado.

  \subsection{Conflictos y alternativas de horario}

  STEM RRL marca las franjas reservadas por otros usuarios y exige que los
  bloques de una misma selección sean contiguos \cite{gilVera2024}.
  Biblio-Smart permite consultar los atributos de distintos asientos libres
  para elegir el que mejor se ajuste a las necesidades del usuario
  \cite{garciaGranja2022}. ITC 2019 distingue las restricciones obligatorias
  de las preferencias al evaluar un horario \cite{muller2025}. CLAB podría
  aplicar esa distinción a cada solicitud. Un solapamiento con una reserva
  confirmada del mismo laboratorio impediría aceptarla. La cercanía a la hora
  solicitada serviría para ordenar las alternativas disponibles. La regla
  contra solapamientos también tendría que verificarse con solicitudes
  simultáneas cuando se implemente la reserva.

  \subsection{Confirmación y uso efectivo del recurso}

  STEM RRL genera una credencial asociada a la reserva para acceder al equipo
  remoto durante la franja asignada \cite{gilVera2024}. Biblio-Smart envía una
  confirmación por correo, detecta la presencia del usuario en el asiento
  mediante el teléfono y libera el puesto después del periodo de ausencia
  configurado \cite{garciaGranja2022}. ITC 2019 genera una asignación que indica
  dónde y cuándo debe realizarse cada clase \cite{muller2025}. CLAB podría
  tomar de Biblio-Smart la distinción entre una reserva confirmada y el uso
  efectivo del espacio. Si el alcance incluye un registro de asistencia, esa
  distinción permitiría identificar qué reservas llegaron a utilizarse.

  \subsection{Bloqueos y cambios de disponibilidad}

  STEM RRL permite habilitar o deshabilitar recursos, fijar sus periodos de
  disponibilidad y programar mantenimiento \cite{gilVera2024}. Biblio-Smart
  actualiza el estado de los asientos según la reserva y la presencia del
  usuario. También simula cómo una regla de distanciamiento cambia los
  puestos disponibles \cite{garciaGranja2022}. ITC 2019 admite periodos en
  los que un aula no puede asignarse \cite{muller2025}. CLAB podría combinar
  el caso de mantenimiento de STEM RRL con los periodos no disponibles de
  ITC 2019 para bloquear un laboratorio durante una franja concreta. Cuando
  un bloqueo afecte una reserva confirmada, el flujo tendría que contemplar
  cómo avisar al responsable y gestionar la reprogramación.

  \subsection{Arquitectura e integración}

  STEM RRL conecta un sistema de reservas independiente con la plataforma
  académica y los equipos de acceso remoto \cite{gilVera2024}. Biblio-Smart
  combina una aplicación móvil, un servicio web, dispositivos de detección
  y un modelo digital del edificio \cite{garciaGranja2022}. ITC 2019 usa
  datos exportados de UniTime para comparar algoritmos que generan horarios
  \cite{muller2025}. En el esqueleto de CLAB, Reservas consulta a
  Laboratorios mediante REST con JSON. El proyecto puede tomar de STEM RRL
  la separación entre la gestión del recurso y la de sus reservas, y de
  Biblio-Smart la idea de presentar el estado del recurso de forma clara
  en la interfaz.

  \subsection{Evaluación del sistema}

  STEM RRL recogió comentarios de 28 estudiantes después de probar el
  laboratorio remoto en tres prácticas, con preguntas sobre la reserva y el
  acceso \cite{gilVera2024}. Biblio-Smart obtuvo 242 respuestas completas y
  voluntarias sobre la calidad percibida del prototipo y la información de
  los asientos \cite{garciaGranja2022}. ITC 2019 comparó algoritmos con 30
  casos de planificación universitaria, entre ellos uno de Yachay Tech,
  Ecuador, y examinó los resultados de asignación y los tiempos de cómputo
  \cite{muller2025}. CLAB podría adoptar dos formas de evaluación. Con usuarios
  de la UTEQ, comprobar si pueden encontrar y solicitar una franja. Con
  pruebas técnicas, verificar que las reglas se cumplan en reservas
  simultáneas, recurrentes y afectadas por bloqueos.
}

% Punto 6 de la guía, dentro de la sección de requisitos del equipo.
\newcommand{\CarlosRequisitosNoFuncionales}{%
  \subsection{Requisitos no funcionales}

  Estos requisitos describen cómo deberá comportarse la versión distribuida de
  CLAB cuando estén implementadas las funciones que se evalúan. Se clasifican
  según el modelo de calidad de
  \href{https://www.iso.org/standard/78176.html}{ISO/IEC 25010:2023}.
  Los 50 usuarios concurrentes, los 2 segundos y el 99\,\% de disponibilidad
  proceden del informe previo de CLAB. La configuración de las pruebas y los
  demás umbrales son propuestas que el equipo deberá validar con el
  levantamiento de información.

  \paragraph{RNF-01. Tiempo de respuesta (eficiencia de desempeño).}
  Durante 10 minutos de carga estable con 50 usuarios virtuales concurrentes,
  al menos el 95\,\% de las consultas válidas de disponibilidad deberá
  devolver el resultado correcto en menos de 2 segundos. La medición se
  realizará desde el cliente de prueba y el cálculo incluirá todas las
  consultas emitidas, contando las respuestas erróneas y los tiempos de
  espera agotados como incumplimientos. Para repetir la prueba se registrarán
  la configuración del entorno, los datos de la agenda y el perfil de carga.
  El uso de percentiles bajo distintos niveles de concurrencia se apoya en
  \cite{pantelic2026Carga}.

  \paragraph{RNF-02. Disponibilidad de la agenda (fiabilidad).}
  Durante dos semanas, dentro del horario académico acordado con la UTEQ, al
  menos el 99\,\% de las comprobaciones programadas de disponibilidad deberá
  devolver un resultado correcto antes de agotar un tiempo de espera de
  5 segundos. Una sonda externa al servicio ejecutará una consulta por minuto
  con una cuenta de prueba autorizada. La tasa se calculará dividiendo las
  comprobaciones exitosas por las programadas en la ventana de observación.
  Se toma como referencia el uso de comprobaciones externas para medir la
  disponibilidad \cite{mulahuwaish2022Disponibilidad}.

  \paragraph{RNF-03. Autorización de operaciones (seguridad).}
  En el 100\,\% de los casos de prueba de confirmación, cancelación y bloqueo,
  la operación deberá quedar sin ejecutar si el usuario carece del permiso,
  la sesión se ha revocado o el servicio de Identidad no puede verificar la
  autorización. Se probará cada operación en esas tres condiciones y se
  comprobará, además de la respuesta, que el estado de la agenda se conserve.
  Los estudios sobre control de acceso y revocación sustentan la verificación
  de estos escenarios
  \cite{venckauskas2023Autorizacion,rasori2024Revocacion}.

  \paragraph{RNF-04. Interfaz en escritorio y móvil (capacidad de interacción).}
  Una persona con permisos deberá poder completar la consulta de disponibilidad
  y el envío de una solicitud en las seis combinaciones de Chrome, Firefox y
  Edge con anchos de pantalla de 360 y 1280 píxeles. En cada combinación
  deberán poder utilizarse los controles necesarios de ambos flujos sin
  desplazamiento horizontal de la página. La prueba registrará la versión
  de cada navegador y el resultado de los dos flujos para cada ancho.
  La evaluación en distintas configuraciones se apoya en el estudio de
  incompatibilidades de presentación \cite{watanabe2023Interfaz}.

  \paragraph{RNF-05. Instalación institucional (flexibilidad).}
  El 100\,\% de los servicios incluidos en la versión evaluada deberá poder
  desplegarse en un entorno limpio, con los prerrequisitos documentados,
  equivalente al entorno institucional. Una vez instalados los servicios,
  la prueba deberá completar una consulta de disponibilidad y confirmar una
  reserva con el acceso a servicios de nube pública deshabilitado.
  El procedimiento toma como referencia la instalación local reproducible
  de una aplicación de microservicios \cite{simoes2026Despliegue}.

  \paragraph{RNF-06. Continuidad ante la caída de Avisos (fiabilidad).}
  Con Información y Avisos detenido y las dependencias de la agenda operativas,
  las 10 solicitudes válidas de la prueba, cada una para una franja libre
  diferente, deberán confirmarse y quedar guardadas en Reservas y Agenda.
  Se comprobará la confirmación y la persistencia de cada solicitud antes
  de restablecer el servicio detenido. Este escenario aplica la evaluación
  mediante fallos inducidos a una dependencia concreta de CLAB
  \cite{alSaidAhmad2024Fallos}.

  \paragraph{RNF-07. Trazabilidad de cambios de agenda (seguridad).}
  El 100\,\% de las confirmaciones, cancelaciones y bloqueos exitosos del
  conjunto de prueba deberá tener un registro de auditoría en el servicio
  de origen. El registro deberá contener el identificador de la operación,
  el usuario, la acción y la fecha y hora. Se cotejará cada cambio con su
  registro tanto con la vista central de auditoría activa como durante su
  caída. Los datos del registro se apoyan en el modelo de trazabilidad de
  eventos de \cite{lopezPimentel2022Auditoria}.
}

% Punto 12 de la guía.
\newcommand{\CarlosRiesgosTecnicos}{%
  \section{Análisis de riesgos}

  El \href{https://github.com/ereinosov/app-dist-form-ac-01/blob/fccfe3f42f0cfa3c405825c4c642b5a5cbfa052b/FORM-AC-01/main.pdf}{informe de FORM-AC-01}
  analiza la caída de Amazon S3 en la región US-EAST-1 del 28 de febrero de
  2017. El \href{https://aws.amazon.com/message/41926/}{informe oficial de AWS}
  documenta tres lecciones que se aplican a CLAB: limitar el alcance de las
  operaciones de mantenimiento, mantener disponible un canal para comunicar
  incidentes y ensayar la recuperación de los componentes críticos. Los
  primeros tres riesgos derivan de ese caso. Los dos restantes surgen de la
  arquitectura propuesta en ADR-001.

  La valoración es cualitativa y corresponde al escenario previo a aplicar
  las mitigaciones. La probabilidad se considera media cuando el escenario
  puede surgir en tareas habituales de mantenimiento, uso o recuperación,
  y baja cuando requiere un fallo excepcional. El impacto es alto si se
  interrumpe la agenda o se compromete la integridad de sus datos. Se considera
  medio si se afecta el diagnóstico o la información secundaria y la agenda
  conserva sus datos. Las verificaciones se proponen para un entorno de
  pruebas con datos de prueba.

  \paragraph{R-01. Bloqueo de la agenda por una operación de mantenimiento.}
  En S3, un parámetro incorrecto permitió retirar más capacidad de la
  prevista. La tabla~2 de FORM-AC-01 traslada esta lección a CLAB: una
  operación de mantenimiento que retenga bloqueos en PostgreSQL podría
  agotar las conexiones de Reservas y Agenda. \textbf{Probabilidad media}
  por la ejecución prevista de tareas de mantenimiento.
  \textbf{Impacto alto} porque se podrían interrumpir las consultas y
  confirmaciones. La mitigación consiste en separar las credenciales de
  la aplicación y de mantenimiento, limitar los permisos y el alcance de
  cada operación, y acotar las esperas por bloqueos y conexiones. Para
  verificarla se mantendrá un bloqueo controlado en la base de pruebas.
  Se comprobará que la API devuelva un error dentro del límite configurado,
  que la solicitud afectada no quede confirmada y que las consultas se
  recuperen al liberar el bloqueo.

  \paragraph{R-02. Pérdida del canal de estado durante una caída.}
  La consola que permitía actualizar el panel de estado de AWS dependía
  del propio S3. En CLAB, si la detección y la comunicación de incidentes
  dependieran de la misma base que usa la agenda, el equipo podría perder
  ambas al mismo tiempo. \textbf{Probabilidad media} por la posibilidad de
  compartir infraestructura. \textbf{Impacto medio} porque retrasaría el
  diagnóstico y el aviso a los usuarios. La mitigación es ejecutar la sonda
  de RNF-02 desde fuera de la aplicación y acordar un canal institucional
  de aviso que pueda usarse durante la caída. Se verificará deteniendo
  PostgreSQL en el entorno de pruebas y comprobando que la sonda registre
  la indisponibilidad y que el aviso pueda emitirse por el canal alterno.

  \paragraph{R-03. Recuperación lenta o incompleta de la agenda.}
  En S3, el reinicio completo y la validación de metadatos tardaron más
  de lo previsto. En CLAB, una recuperación completa de PostgreSQL podría
  tardar más de lo esperado o devolver una agenda que requiera comprobaciones
  adicionales. \textbf{Probabilidad baja} porque el escenario requiere
  una interrupción que obligue a recuperar el almacén completo.
  \textbf{Impacto alto} por la posible interrupción prolongada de las
  confirmaciones. La mitigación es ensayar el procedimiento de recuperación
  y verificar la agenda antes de habilitar de nuevo las escrituras. Se
  restaurará una copia conocida, se medirá el tiempo hasta recuperar una
  consulta y una confirmación correctas, y se cotejarán las reservas y
  bloqueos recuperados con los datos esperados de esa copia.

  \paragraph{R-04. Confirmación de reservas superpuestas.}
  Si dos solicitudes concurrentes validan la misma franja antes de que
  se guarde la primera, ambas podrían confirmarse. \textbf{Probabilidad
  media} porque varios usuarios pueden solicitar el mismo laboratorio.
  \textbf{Impacto alto} porque se asignaría un espacio a actividades
  incompatibles. Según ADR-001, Reservas y Agenda debe serializar las
  modificaciones por laboratorio y verificar los solapamientos dentro de
  la misma transacción que guarda el cambio. Se enviarán dos solicitudes
  válidas para el mismo laboratorio y franja, y se comprobará que quede
  una sola reserva confirmada. Se repetirá la prueba con una reserva frente
  a un bloqueo para comprobar que no queden ambos vigentes.

  \paragraph{R-05. Avisos omitidos o duplicados al recuperar el servicio.}
  Si Información y Avisos se detiene mientras consulta los cambios de la
  agenda, podría omitir un evento pendiente o procesarlo dos veces al
  reiniciar. \textbf{Probabilidad media} porque una interrupción puede
  coincidir con el procesamiento. \textbf{Impacto medio} porque los
  usuarios podrían recibir avisos repetidos o consultar una vista de
  auditoría incompleta. La mitigación propuesta en ADR-001 consiste en
  conservar un registro durable en el servicio de origen, guardar el
  avance del consumidor y deduplicar por origen e identificador. Se
  detendrá Información y Avisos, se confirmarán reservas durante el corte
  y, al reanudarlo, se comprobará que cada evento pendiente se refleje
  una sola vez. La prueba complementa RNF-06 y RNF-07.
}
